\subsection{角域。}

我们先不作任何进一步的假设，而设 E 是有序域 A 上的一个定向平面。我们说 E 的三条半直线 \(D_{0}\) 、 \(D_{1}\) 、 \(D_{2}\) （以 0 为原点）构成正向序列，如果对 \(x_{i} \in D_{i}\) ， \(x_{i} \neq 0\) （i = 0, 1, 2），二重向量 \(x_{0} \wedge x_{1}\) 、 \(x_{1} \wedge x_{2}\) 、 \(x_{2} \wedge x_{0}\) 中至少有两个 \textgreater0；此时，序列 \(D_{1}\) 、 \(D_{2}\) 、 \(D_{0}\) 与 \(D_{2}\) 、 \(D_{0}\) 、 \(D_{1}\) 也是正向的。显然，构成正向序列的三条半直线两两不同。给定 E 的两条半直线 \(D_{1}\) 、 \(D_{2}\) ，以 \(D_{1}\) 为始边、 \(D_{2}\) 为终边的开（相应地，闭）角域是指使序列 \(D_{1}\) 、D、 \(D_{2}\) 为正向（相应地，使 \(D = D_{1}\) 或 \(D = D_{2}\) ，或使序列 \(D_{1}\) 、D、 \(D_{2}\) 为正向）的半直线 D 所成的集（或者，滥用术语，这些半直线的并）。

命题 11。—— 设 E 是有序域 A 上的一个定向平面，D₀ 为 E 的一条半直线，G 为 E 中异于 D₀ 的半直线所成的集。关系

“\(D_{1} = D_{2}\)，或序列 \(D_{0}\) 、 \(D_{1}\) 、 \(D_{2}\) 为正向”

（定义于 G 的元素 \(D_{1}\) 、 \(D_{2}\) 之间）是 G 上的一个全序关系。

事实上，全序关系的各条公理除传递性外均可平凡地验证。设 \(D_{1}\) 、 \(D_{2}\) 、 \(D_{3}\) 为三条半直线，使得序列 \(D_{0}\) 、 \(D_{1}\) 、 \(D_{2}\) 与 \(D_{0}\) 、 \(D_{2}\) 、 \(D_{3}\) 皆为正向；我们要证明序列 \(D_{0}\) 、 \(D_{1}\) 、 \(D_{3}\) 是正向的。为此，取向量 \(x_{i} \neq 0\) 属于 \(D_{i} (i = 0, 1, 2, 3)\) ，选取二重向量 e 使 \textgreater{} 0，并令 \(x_{i} \wedge x_{j} = a_{ij} e (a_{ij} \in A)\) 。写 \(e = x_{0} \wedge y (y \in E)\) ，并取 \((x_{0}, y)\) 作为 E 的基，即可毫不费力地验证关系式

\[
a _ {0 1} a _ {2 3} + a _ {0 2} a _ {3 1} + a _ {0 3} a _ {1 2} = 0.\tag{22}
\]

如此，若 \(a_{01} \leqslant 0\) ，则有 \(a_{12} > 0\) 与 \(a_{20} > 0\) （因第一个序列为正向），进而有 \(a_{23} > 0\) 与 \(a_{30} > 0\) （因 \(a_{02} < 0\) 且第二个序列为正向），从而 \(a_{13} > 0\) （由 (22)）；因此在这种情形下序列 \((\mathrm{D}_0, \mathrm{D}_1, \mathrm{D}_3)\) 是正向的。现设 \(a_{01} > 0\) 。若 \(a_{30} \leqslant 0\) ，则有 \(a_{02} > 0\) 与 \(a_{23} > 0\) （因第二个序列为正向），进而有 \(a_{12} > 0\) （因 \(a_{20} < 0\) 且第一个序列为正向），故 \(a_{13}>0\) （由 (22)），于是序列 \((\mathrm{D}_{0}, \mathrm{D}_{1}, \mathrm{D}_{3})\) 是正向的。最后，若 \(a_{01}>0\) 且 \(a_{30}>0\) ，情形显然相同。证毕。

推论。—— 设 \(D_{1}\) 与 \(D_{2}\) 为 E 的两条不同的半直线。对 E 的每条半直线 \(D_{0}\) ，当序列 \(D_{0}\) 、 \(D_{1}\) 、 \(D_{2}\) 为正向时，满足 \(D_{1} < D < D_{2}\) （对由 \(D_{0}\) 所定义的全序关系而言）的 E 的半直线 D 所成的集，等于以 \(D_{1}\) 为始边、 \(D_{2}\) 为终边的开角域。

事实上，给定一条半直线 \(D_{3}\) ，我们必须证明关系“序列 \(D_{1}\) 、 \(D_{3}\) 、 \(D_{2}\) 为正向”与“序列 \(D_{0}\) 、 \(D_{1}\) 、 \(D_{3}\) 与 \(D_{0}\) 、 \(D_{3}\) 、 \(D_{2}\) 皆为正向”等价。为简便起见，用 (ijk) 表示关系“序列 ( \(D_{i}\) , \(D_{j}\) , \(D_{k}\) ) 为正向”。由命题 11，(132) 与 (120) 的合取蕴含 (130)；同样，(201) 与 (213) 的合取蕴含 (203)；由此得断言的一半。反之，设 (012)、(013) 与 (032)；由于 (312) 与 (320) 的合取蕴含 (310)（命题 11），而 (310) 与 (013) 不相容，故 (312) 为假；这就证明了 (132)，从而完成证明。

由上述推论，以 \(D_{1}\) 为始边、 \(D_{2}\) 为终边的开（相应地，闭）角域记作 \(\left[\begin{matrix}\mathrm{D}_{1},\mathrm{D}_{2}\end{matrix}\right]\) （相应地， \(\left[\begin{matrix}\mathrm{D}_{1},\mathrm{D}_{2}\end{matrix}\right]\) ），这里应理解为：它们是对于由任一半直线 \(D_{0}\) （序列 \(D_{0},D_{1},D_{2}\) 为正向）所定义的序结构而言的区间。

命题 12。—— 设 A 为极大有序域，E 为 A 上的一个定向平面，D₀ 为 E 的一条半直线，G 为 E 中异于 D₀ 的半直线所成的集。全序集 A 与 G（命题 11）同构。

事实上，设 \((x, y)\) 为 E 的一个基，使得 \(x \in -\mathrm{D}_{0}\) 且二重向量 \(x \wedge y\) 为 \(> 0\) 。对 A 的每个元素 \(t\) ，让它对应半直线 \(f(t)\) ，即向量 \((1 - t^{2})x + 2ty\) 所属的那条半直线。显然 \(f(\mathrm{A}) \subset \mathrm{G}\) 。我们证明 \(f\) 是严格递增的。事实上，按定义，要使序列 \(\mathrm{D}_{0}\) 、 \(f(t)\) 、 \(f(t')\) （ \(t, t'\) 属于 A）为正向，必须且只需下列元素

\[
- 2 t, \quad (1 - t ^ {2}) 2 t ^ {\prime} - (1 - t ^ {\prime 2}) 2 t, \quad 2 t ^ {\prime}
\]

中至少有两个 \textgreater0。而第二个元素等于 \(2(t' - t)(1 + tt')\) 。因此，若 \(t < t'\) ，则或者 \(tt' \geqslant 0\) ，于是 \(t' > 0\) 或 -t \textgreater{} 0；或者 \(tt' < 0\) ，于是 -t \textgreater{} 0 且 \(t' > 0\) ；在任何情形下 \(D_{0}\) 、 \(f(t)\) 、 \(f(t')\) 都为正向。由于 A 是全序的，f 是 A 到有序集 \(f(A)\) 上的一个同构（《集合论》，第 III 章，§ 1，第 14 小节，命题 13）。

剩下要证明 \(f\) 是满射。为此，考虑正型 \(\Phi\) （于 \(E\) 上），使 \((x, y)\) 为关于 \(\Phi\) 的标准正交基。对每条半直线 \(D \in G\) ，存在且仅存在一个角 \(\varphi\) 使 \(2\varphi = (-D_0, D)\) （ \(n^o 1\) ，命题 3）；由于 \((-D_0, D)\) 不是平角，\(\varphi\) 不是直角，因而 \(tg\varphi\) 取有限值。于是，由公式 (16)（ \(n^o 3\) ），我们有 \(D = f(tg\varphi)\) 。证毕。

习题。—— 1) 采用第 1 小节的记号，令 Q(x) = Φ(x, x)；向量空间 E 典范地等同于 C⁻(Q)（§ 9，第 1 小节）。对每个 z ∈ C⁺(Q) 与每个 x ∈ E，有 zx ∈ E；证明 x ↦ zx 是 A(Φ) 的一个元素，且 z ↦ sz 是 C⁺(Q) 到代数 A(Φ) 上的一个同构。

\begin{enumerate}
\def\labelenumi{\arabic{enumi})}
\setcounter{enumi}{1}
\tightlist
\item
  假设与记号同第 1 小节及习题 1。
\end{enumerate}

\begin{enumerate}
\def\labelenumi{\alph{enumi})}
\item
  设 C 是 \(x \in E\) 中满足 \(\Phi(x, x) = 1\) 的那些 x 所成的集（“单位圆”），并设 \(\mathfrak{D}\) 为与 C 的交非空的直线 D 所成的集（因而该交由 E 的两个相反的元素组成）。称偶 \(\Delta = (D, z)\) （由一条直线 \(D \in \mathfrak{D}\) 与点 \(z \in D \cap C\) 之一组成）为带基点直线。证明：若 \(\Delta_1 = (D_1, z_1)\) 、 \(\Delta_2 = (D_2, z_2)\) 是两条带基点直线，则存在且仅存在一个旋转 \(u\) 使 \(u(z_1) = z_2\) （因而 \(u(D_1) = D_2\) ），把这写成 \(u(\Delta_1) = \Delta_2\) 。在带基点直线的偶 \((\Delta_1, \Delta_2)\) 所成的集上，关系“存在一个旋转 \(u\) 使 \(u(\Delta_1) = \Delta_1'\) 且 \(u(\Delta_2) = \Delta_2'\) ”是一个等价关系。在这一关系下带基点直线的等价类所成的集 \(\mathfrak{A}_1\) 称为带基点直线的角所成的集，而带基点直线偶 \((\Delta_1, \Delta_2)\) 所属的等价类称为该偶的角，记作 \((\widehat{\Delta_1, \Delta_2})\) ；关系 \((\widehat{\Delta_1, \Delta_2}) = (\widehat{\Delta_1'}, \widehat{\Delta_2'})\) 等价于 \((\widehat{\Delta_1, \Delta_1'}) = (\widehat{\Delta_2, \Delta_2'})\) ，且把 \(\Delta_1\) 映为 \(\Delta_2\) 的旋转称为角为 \(\theta = (\Delta_1, \Delta_2)\) 的旋转，记作 \(h_1(\theta)\) ； \(h_1\) 是 \(\mathfrak{A}_1\) 到 \(O^+\) 上的一个双射，且 \(\mathfrak{A}_1\) 借助 \(h_1^{-1}\) 获得 \(O^+\) 的交换群结构，并把如此定义的群 \(\mathfrak{A}_1\) 写成加法记号；又称 \(\mathfrak{A}_1\) 中由带基点直线 \((D, z)\) 与“相反的”带基点直线 ( \(D, -z\) ) 所成之偶的角为平角，它对应于旋转 \(x \to -x\) 。
\item
  在非迷向直线所成的集 \(\mathfrak{D}_{0} \supset \mathfrak{D}\) 中，如同第 3 小节那样定义直线间的角的概念、群 \(\mathfrak{A}_{0}\) （利用第 1 小节命题 1 的推论 3）、 \(\mathfrak{A}_{0}\) 中的直角，以及典范双射 \(h\) （从 \(\mathfrak{A}_{0}\) 到 \(S^{+}/H\) 上）。用命题 3 之推论的记号，令 \(\bar{d} = h_{1}^{-1} \circ d \circ h\) 与 \(\bar{i} = h^{-1} \circ i \circ h_{1}\) ，使得 \(\bar{i}\) 是 \(\mathfrak{A}_{1}\) 到 \(\mathfrak{A}_{0}\) 中的同态，而 \(\bar{d}\) 是 \(\mathfrak{A}_{0}\) 到 \(\mathfrak{A}_{1}\) 中的同态； \(\bar{d}\) 是双射，且 \(\bar{i}\) 的核由 0 与平角组成；此外，我们有 \(\bar{d}(\bar{i}(\theta)) = 2\theta\) （对 \(\theta \in \mathfrak{A}_{1}\) ），以及 \(\bar{i}(\bar{d}(\varphi)) = 2\varphi\) （对 \(\varphi \in \mathfrak{A}_{0}\) ）。要使 \(\bar{i}\) 为满射（换言之，要使 \(\mathfrak{D}\) 成为全体非迷向直线所成的集），必须且只需域 A 是毕达哥拉斯域（第 VI 章，§ 2，习题 8 d)），且存在关于 \(\Phi\) 的标准正交基；这等价于说 \(\Phi(x, x)\) 对每个 \(x \in E\) 都是平方元。
\end{enumerate}

\begin{enumerate}
\def\labelenumi{\arabic{enumi})}
\setcounter{enumi}{2}
\tightlist
\item
  假设与记号同习题 2。
\end{enumerate}

\begin{enumerate}
\def\labelenumi{\alph{enumi})}
\tightlist
\item
  设 \(s\) 是关于一条非迷向直线 D 的对称（变换）（ \(\S 6\) ，第 4 小节）。对每条带基点直线 \(\Delta_{1} = (D_{1}, z_{1})\) ，令
\end{enumerate}

\[
\Delta_ {2} = s (\Delta_ {1}) = (s (\mathrm{D} _ {1}), s (z _ {1})),
\]

并令 \(\varphi = (\widehat{\mathrm{D}_1, \mathrm{D}})\) ；证明有 \((\widehat{\Delta_1, \Delta_2}) = \bar{d}(\varphi)\) 。

\begin{enumerate}
\def\labelenumi{\alph{enumi})}
\setcounter{enumi}{1}
\item
  证明每个行列式为 -1 的正交变换都是关于一条非迷向直线的对称（变换） \(s\) （参见 § 6，习题 15 e)）；对每个旋转 \(u\) ，有 \(sus^{-1} = u^{-1}\) 。
\item
  若 \(x, y\) 是 E 的任意两点，满足 \(\Phi(x, x) = \Phi(y, y) \neq 0\) ，则存在且仅存在一个关于非迷向直线的对称（变换），它把 \(x\) 变为 \(y\) 。
\item
  设 \(s, s'\) 为关于两条非迷向直线 D、D' 的对称（变换），并令 \(\varphi = (\widehat{\mathrm{D},\mathrm{D}'})\) ；要使 \(s's\) 是角为 \(\theta\) 的旋转，必须且只需 \(\bar{d}(\varphi) = \theta\) 。
\item
  证明正交群 \(\mathbf{O}(\mathrm{Q})\) 的换位子群是 \(\mathfrak{A}_{1}\) 在同态 \(\theta \to h_{1}(2\theta)\) 下的像（ \(\S 6\) ，习题 \(17a\) ）；要使这个像等于 \(\mathrm{O}^{+}\) ，必须且只需同态 \(\bar{i}\) 是满射（习题 \(2b\) ）。
\end{enumerate}

\begin{enumerate}
\def\labelenumi{\arabic{enumi})}
\setcounter{enumi}{3}
\item
  假设与记号同习题 2。设 \(a\) 、 \(b\) 为圆 C 的两点， \(\Delta_{a}\) 、 \(\Delta_{b}\) 分别是过 \(a\) 与 \(b\) （以及过 0）的带基点直线，并令 \(\theta = (\widehat{\Delta_{a}, \Delta_{b}})\) 。对每个 \(x \in E\) ，x 异于 \(a\) 与 \(b\) ，令 D \(_{xa}\) （相应地，D \(_{xb}\) ）为过 \(a\) 与 \(x\) （相应地，过 \(b\) 与 \(x\) ）的仿射直线，并令 D \(_{xa}'\) （相应地，D \(_{xb}'\) ）为 D \(_{xa}\) （相应地，D \(_{xb}\) ）的方向；证明：要使 \(x \in C\) （ \(x\) 异于 \(a\) 与 \(b\) ），必须且只需角 \(\varphi = (\widehat{\mathrm{D}_{xa}', \mathrm{D}_{xb}'})\) 满足 \(\bar{d}(\varphi) = \theta\) （利用习题 3）。当 \(x = a\) 或 \(x = b\) 时，这一结果如何变化（参见 § 6，习题 25) \(a\) ))？
\item
  采用第 1、2 小节及习题 2 的记号，设 \(\Phi\) 的指数为 1，换言之，\(\delta\) 是 A 中的平方元 \(\gamma^{2}\) ；用 D \(_{1}\) 、D \(_{2}\) 表示 E 的迷向直线，它们分别包含向量 \(e_{1}-\frac{1}{\gamma}e_{2}\) 与 \(e_{1}+\frac{1}{\gamma}e_{2}\) 。
\end{enumerate}

\begin{enumerate}
\def\labelenumi{\alph{enumi})}
\item
  证明旋转群 \(O^{+}\) 同构于域 A 的乘法群 A*。
\item
  对每个角 \(\theta \in \mathfrak{A}_1\) ，令 \(e_w(\theta) = c_w(h_1(\theta)) + \gamma s_w(h_1(\theta))\) 。证明 \(\theta \to e_w(\theta)\) 是 \(\mathfrak{A}_1\) 到群 \(A^*\) 上的一个同构。
\item
  设 D 与 D′ 为任意两条直线，并令 \(\varphi = (\widehat{\mathbf{D},\mathbf{D}^{\prime}})\) 。证明交比 \(\left[ \begin{array}{cc}\mathrm{D}_{1} & \mathrm{D}_{2}\\ \mathrm{D}^{\prime} & \mathrm{D} \end{array} \right]\) （第 II 章，第 2 版，附录 III，习题 5）等于 \(e_w(\overline{d} (\varphi))\) （“拉盖尔公式”）。（注意 D1 与 D2 在任何正相似变换下不变，并利用第 II 章第 2 版附录 III 的习题 4 c)，化归到 D = Ae1 的情形。）
\end{enumerate}

\begin{enumerate}
\def\labelenumi{\arabic{enumi})}
\setcounter{enumi}{5}
\tightlist
\item
  设 A 是一个有序域。
\end{enumerate}

\begin{enumerate}
\def\labelenumi{\alph{enumi})}
\item
  设 \(D_{1}\) 、 \(D_{2}\) 为以 0 为原点的两条非迷向半直线。证明存在一个正相似变换 u 使 \(u(\mathrm{D}_{1}) = \mathrm{D}_{2}\) ；任何具有这一性质的其他正相似变换都具有 \(\varrho = su\) 的形式，其中 s 是比为 \textgreater{} 0 的位似。
\item
  在非迷向半直线的偶 \((\mathrm{D}_{1}, \mathrm{D}_{2})\) 所成的集上，关系“存在正相似变换 u 使 \(u(\mathrm{D}_{1}) = \mathrm{D}_{1}'\) 且 \(u(\mathrm{D}_{2}) = \mathrm{D}_{2}'\) ”是一个等价关系。非迷向半直线在这一关系下的等价类所成的集 A 称为（非迷向）半直线的角所成的集，而这类半直线的偶 \((\mathrm{D}_{1}, \mathrm{D}_{2})\) 的等价类称为该偶的角，记作 \((\widehat{\mathrm{D}_{1}, \mathrm{D}_{2}})\) ；若 \(\theta = (\widehat{\mathrm{D}_{1}, \mathrm{D}_{2}})\) ，则称 \(\theta\) 为每个把 \(D_{1}\) 映为 \(D_{2}\) 的正相似变换的角；用 \(h_{2}(\theta)\) 表示这些相似变换关于模 \(H^{+}\) 的类，于是 \(h_{2}\) 是集 A 到群 \(S^{+}/H^{+}\) 上的一个双射；借助 \(h_{2}^{-1}\) 把 \(S^{+}/H^{+}\) 的交换群结构转移到 A 上，并把如此定义的群 A 写成加法记号。试定义一个典范单射 \(\bar{j}\) ，它把带基点直线的角群 \(A_{1}\) （习题 2）映入群 A，使 \(h_{2} \circ \bar{j} \circ h_{2}^{-1}\) 是 \(O^{+}\) 到 \(S^{+}/H^{+}\) 中的典范单射 j。要使 \(\bar{j}\) 为满射，必须且只需同态 \(\bar{i}\) （ \(A_{1}\) 到 \(A_{0}\) 中的）是满射（习题 2 b)）。
\item
  证明在 \(\mathfrak{A}\) 中，方程 \(2\theta = 0\) 当 \(\delta < 0\) 时有两个解，当 \(\delta > 0\) 时有四个解。在第一种情形，该方程的解 \(\varpi \neq 0\) 也称为平角。
\end{enumerate}

¶ 7) 在习题 6 的假设与记号下，设同态 \(\bar{j}\) 是双射；于是可如第 3 小节那样对每个 θ ∈ A 定义 cos θ 与 sin θ。设 T 是 A 中满足 sin θ ≥ 0 的 θ 所成的集。

\begin{enumerate}
\def\labelenumi{\alph{enumi})}
\item
  证明对每个 \(\theta \in T\) ，存在且仅存在一个角 \(\theta' \in T\) 使 \(2\theta' = \theta\) ；令 \(\theta' = \theta/2\) 。
\item
  设 L 是以 Z 为系数、由 T 的元素的形式线性组合构成的 Z-模（第 II 章，§ 1，第 8 小节）；用 \(\xi + \eta\) 与 \(\dot{\cdot} \xi\) 分别表示 L 中的和与相反元。设 N 是 L 中由形如 \(\xi + \eta \dot{\cdot} (\xi + \eta)\) 的元素（对一切元素偶 \((\xi, \eta)\) ，其中 \(\xi + \eta \in T\) ，和取在群 \(\mathfrak{A}\) 中）生成的子模。设 \(f\) 是 L 到 \(\mathfrak{A}\) 的同态，它开拓 T 到 \(\mathfrak{A}\) 中的典范单射， \(g\) 是 L 的自同态，它开拓 T 到自身的映射 \(\theta \to \theta/2\) 。我们有 \(f(N) = \{0\}\) 与 \(g(N) \subset N\) ；通过过渡到商，由 \(f\) 得出一个同态 \(f\) （从 M = L/N 到 \(\mathfrak{A}\) ），由 \(g\) 得出 M 的一个自同态 \(g\) ；令 \(g(\mu) = \mu/2\) ，而若 \(g^m\)
\end{enumerate}

是 g 的 m 次迭代，则 \(g^{m}(\mu) = 2^{-m}\mu\) ；我们有 \(2^{m}(2^{-m}\mu) = \mu\) （对一切 \(\mu \in M\) ）。

\begin{enumerate}
\def\labelenumi{\alph{enumi})}
\setcounter{enumi}{2}
\item
  证明 L 到 M = L/N 上的典范映射 \(\psi\) 在 T 上的限制是单射，从而可借助 \(\psi\) 把 T 等同于 M 的一个子集。证明：若 \(\lambda_1, \ldots, \lambda_m\) 是 T 中 \(\neq 0\) 的元素，则和 \(\lambda_1 + \lambda_2 + \cdots + \lambda_m\) 在 M 中不能为 0（考察元素 \(2^{-m}(\lambda_1 + \cdots + \lambda_m)\) 并对 m 作归纳，注意这些元素属于 T）。
\item
  设 \(M_{+}\) 是 T 的元素（在 M 中）的有限和所成的集；证明 \(\mathbf{M}_{+}\cap(-\mathbf{M}_{+})=\left\{0\right\}\) 且 \(\mathbf{M}=\mathbf{M}_{+}\cup(-\mathbf{M}_{+})\) ，因而 \(M_{+}\) 是 M 上某个全序群结构中 \(\geqslant0\) 的元素所成的集（注意对每个 \(\mu\in M_{+}\) ，存在整数 m 使 \(2^{-m}\mu\in T\) ）；称这个全序群为半直线角的度量群。证明 M 到 A 的同态 f 是满射，且其核是 \(2\varpi\) 的整数倍所成的集，其中 \(\varpi\) 是平角（习题 6 c)）。证明 T 等同于 M 中的区间 \([0,\varpi]\) （对 m 作归纳证明：若 \(\mu,\lambda_{1},\ldots,\lambda_{m}\) 属于 T 且 \(\lambda_{1}+\cdots+\lambda_{m}\leqslant\mu\) ，则 \(\lambda_{1}+\cdots+\lambda_{m}\in T\) ）。证明在 T（因而等同于 M 的一个区间）中，函数 \(\theta\to\cos\theta\) 严格递减。
\item
  全序群 M 为阿基米德群（第 VI 章，§ 1，习题 31），必须且只需域 A 的加法群是阿基米德群。（为看条件必要，注意若 sin θ 关于 A 的子域 Q（第 VI 章，§ 2，习题 1）是无穷小，则对每个整数 n，sin nθ 亦然。为看条件充分，注意若 0 ≤ θ ≤ π/4，则 sin 2θ ≥ √2 sin θ。）
\end{enumerate}

\begin{enumerate}
\def\labelenumi{\arabic{enumi})}
\setcounter{enumi}{7}
\tightlist
\item
  设 E 是有序域 A 上的一个定向平面。设 D'、D'' 为两条不同的半直线；设 \(x'\) （相应地， \(x''\) ）是 D′（相应地，D″）中一个 \(\neq 0\) 的向量；以 D' 为始边、D'' 为终边的角域（开的或闭的）称为凸角域（相应地，凹角域、平角域），如果 \(x' \wedge x'' > 0\) （相应地， \(x' \wedge x'' < 0\) 、 \(x' \wedge x'' = 0\) ）。
\end{enumerate}

\begin{enumerate}
\def\labelenumi{\alph{enumi})}
\item
  要存在向量空间 E 的自同构把开（相应地，闭）角域 \(\Sigma_{1}\) 映入开（相应地，闭）角域 \(\Sigma_{2}\) ，必须且只需 \(\Sigma_{1}\) 与 \(\Sigma_{2}\) 同为凸角域，或同为凹角域，或同为平角域。
\item
  证明有序集 \(\left[D', D''\right]\) 同构于 A 的区间 \(\left[0,1\right]\) （先考虑凸角域的情形，并注意在 A 中任何两个有界闭区间都是同构的有序集）。
\item
  在习题 7 的记号与假设下，定义从集 T 到一个平角域上的典范双射，并证明此映射是序结构之间的同构。
\end{enumerate}

\begin{enumerate}
\def\labelenumi{\arabic{enumi})}
\setcounter{enumi}{8}
\item
  设 A 是一个毕达哥拉斯有序域，E 是 A 上的一个有限维向量空间，Q 是 E 上的一个正的非退化二次型，且 E 关于它具有标准正交基（换言之，对每个 \(x \in E\) ，Q(x) 都是 A 中的平方元）。证明换位子群 \(\Omega(Q)\) （正交群 \(\mathbf{O}(Q)\) 的换位子群）就是旋转群 \(\mathbf{O}^{+}(Q) = \mathbf{S}\mathbf{O}(Q)\) （利用 § 10 习题 3 e) 与 § 6 习题 17 a)）。
\item
  设 A 是一个极大有序域，E 是 A 上的一个有限维向量空间，Q 是 E 上的一个正的非退化二次型。对每个正交变换 \(u \in \mathbf{O}(Q)\) ，证明 E 可分解为两两正交的子空间 P、N、 \(R_i\) （ \(1 \leqslant i \leqslant r\) ）的直和，且具有下列性质： \(1^\circ u(x) = x\) 在 P 中成立， \(u(x) = -x\) 在 N 中成立； \(2^\circ\) 每个 \(R_i\) 的维数为 2，我们有 \(u(R_i) = R_i\) ，且 \(u\) 在 \(R_i\) 上的限制是角为 \(\theta_i\) 的旋转，此角异于 0 且异于平角。此外，对这种类型的两个分解而言，子空间 P 与 N 相同，元素序列 \(\cos \theta_i\) 在相差一个次序的意义下也相同（参见 § 7，第 3 小节，定理 2 的推论 2）。由此推出，每个旋转 \(u \in \mathbf{O}^+(Q)\) 都是一个换位子 \(tst^{-1}s^{-1}\) ，其中 \(s\) 与 \(t\) 属于 \(\mathbf{O}(Q)\) 且 \(s^2 = 1\) （参见习题 3 d)）。
\item
  设 A 是一个极大有序域，L 是 A 上关于数对 \((-1,-1)\) 的四元数体，E 是 L 中由纯四元数组成的（维数为 3 的）子空间（ \(\S 9\) ，习题 15 a)）；于是每个四元数都可唯一地写成 \(s=\alpha.1+\nu\) ，其中 \(\alpha\in A\) 且 \(\nu\in E\) ；我们有 \(\alpha^{2}-\nu^{2}=\rho.1\) ，其中 \(\rho\in A\) 且 \(\rho\geqslant0\) ；令 \(\|\nu\|=\sqrt{\rho}\) 。
\end{enumerate}

\begin{enumerate}
\def\labelenumi{\alph{enumi})}
\item
  设 \(\varphi(s)\) 是 E 中的旋转 \(x \to sxs^{-1}\) ，它关于 E 上的正的非退化二次型 \(x \to \|x\|^2\) 而取（§ 9，第 5 小节，定理 4 及习题 15）。证明：若 \(\nu \neq 0\) ，则包含 \(\nu\) 的直线 D ⊂ E 的向量在 \(\varphi(s)\) 下不变； \(\varphi(s)\) 在正交平面 P = D⁰ 上的限制是角为 \(\theta\) 的旋转，使得（对 P 的一个适当定向）有 \(\operatorname{tg} \frac{\theta}{2} = \| \nu \| / \alpha\) 。（若 (1, i, j, k) 是 L 在 A 上的典范基，则化归到 \(\nu = \beta i\) 、 \(\beta \in A\) 的情形，然后计算 \(sjs^{-1}\) 。）
\item
  证明每个范数为 1 的四元数都可以写成 \(tst^{-1}s^{-1}\) ，其中 \(s \in L\) ， \(t \in E\) （参见习题 10）；并且，若 \(a\) 、 \(b\) 是两个范数为 1 的纯四元数，则存在四元数 \(s\) 使 \(b = sas^{-1}\) 。
\item
  要使两个四元数 \(s = \alpha + \nu\) 、 \(t = \beta + w\) （其中 \(\alpha, \beta\) 属于 A， \(\nu, w\) 属于 E）可交换，必须且只需 \(\nu = \lambda w\) （ \(\lambda \in A\) ）；要使 \(st = -ts\) ，必须且只需 \(\alpha = \beta = 0\) 且 E 中的向量 \(\nu\) 与 \(w\) 正交（参见 § 3，习题 10）。
\end{enumerate}

\begin{enumerate}
\def\labelenumi{\arabic{enumi})}
\setcounter{enumi}{11}
\tightlist
\item
  设 A 是一个极大有序域，L 是 A 上的一个 n 维欧几里得空间，其度量型 \(\Phi\) 是正的非退化型；用 \(d(x, y)\) 表示 L 的两点之间的距离 \(\sqrt{\Phi(x - y, x - y)}\) （ \(\S 7\) ，第 1 小节，注）。对每一点 \(c \in L\) 与 A 的每个元素 \(\rho > 0\) ，以 c 为中心、以 \(\rho\) 为半径的球面是指 \(x \in L\) 中满足 \(d(x, c) = \rho\) 的点所成的集；因此球面是 L 中有中心的非退化仿射二次曲面（ \(\S 6\) ，习题 25）。
\end{enumerate}

\begin{enumerate}
\def\labelenumi{\alph{enumi})}
\item
  证明：在以 c 为中心的球面 S 的每一点 x 处，S 的切超平面（§ 6，习题 25）与过 c 和 x 的直线垂直（§ 6，习题 22）。
\item
  设 S 是以 c 为中心、以 \(\rho\) 为半径的球面，a 是 L 的一点，D 是过 a 且与 S 交于两个不同点 \(x_{1}\) 、 \(x_{2}\) （相应地，与 S 相切于一点 x）的直线。证明有
\end{enumerate}

\[
\Phi (x _ {1} - a, x _ {2} - a) = (d (a, c)) ^ {2} - \rho^ {2} \quad (\text { resp. } (d (x, a)) ^ {2} = (d (a, c)) ^ {2} - \rho^ {2}).
\]

A 的元素 \((d(a, c))^{2} - \rho^{2}\) 称为 a 关于 S 的幂。

\begin{enumerate}
\def\labelenumi{\alph{enumi})}
\setcounter{enumi}{2}
\item
  设 \(S_{1}\) 、 \(S_{2}\) 为两个不同心的球面；证明 L 中关于 \(S_{1}\) 与 \(S_{2}\) 的幂相等的点所成的集是一个超平面，它垂直于过这两个球面中心的直线，且包含它们的交 \(S_{1} \cap S_{2}\) ；称这个超平面为 \(S_{1}\) 与 \(S_{2}\) 的根超平面。
\item
  设 \(S_{1}\) 、 \(S_{2}\) 为两个球面，其中心分别为 \(c_{1}\) 、 \(c_{2}\) ，半径分别为 \(\rho_{1}\) 、 \(\rho_{2}\) 。证明下列性质等价：
\end{enumerate}

α) 交 \(S_{1} \cap S_{2}\) 非空，且对这个交的每一点， \(S_{1}\) 与 \(S_{2}\) 在该点的切超平面互相垂直。

\(\beta_{1}\) \(c_{1}\) 关于 \(S_{2}\) 的幂是 \(\rho_{1}^{2}\) 。

\(\beta_{2})\) \(c_{2}\) 关于 \(\mathbf{S}_{1}\) 的幂是 \(\rho_{2}^{2}\) 。

\(\gamma_{1})\) \(\mathbf{S}_{1}\) 与 \(\mathbf{S}_{2}\) 的根超平面是 \(c_{1}\) 关于 \(\mathbf{S}_{2}\) 的极超平面（§ 6，习题 25）。

\(\gamma_{2})\) \(\mathbf{S}_{1}\) 与 \(\mathbf{S}_{2}\) 的根超平面是 \(c_{2}\) 关于 \(\mathbf{S}_{1}\) 的极超平面。

当这些性质成立时，称球面 \(S_{1}\) 、 \(S_{2}\) 是正交的。

\begin{enumerate}
\def\labelenumi{\alph{enumi})}
\setcounter{enumi}{4}
\tightlist
\item
  设 \(S_{1}\) 、 \(S_{2}\) 是两个正交的球面， \(c_{1}\) 、 \(c_{2}\) 是它们的中心。证明：若 \(\varpi_{1}\) 、 \(\varpi_{2}\) 是一点 x 关于 \(S_{1}\) 、 \(S_{2}\) 的幂，则有 \(\varpi_{1} + \varpi_{2} = 2\Phi(x - c_{1}, x - c_{2})\) 。逆之亦真。
\end{enumerate}

\begin{enumerate}
\def\labelenumi{\arabic{enumi})}
\setcounter{enumi}{12}
\tightlist
\item
  假设与记号同习题 12。给定点 \(c \in \mathrm{L}\) 与一个元素 \(\alpha \neq 0\) （属于 \(\mathbf{A}\) ），我们把以 \(c\) 为极点、以 \(\alpha\) 为幂的反演理解为这样一个对合置换 \(u\) （集 \(\mathrm{L} - \{c\}\) 的）：对一切 \(x \in \mathrm{L} - \{c\}\) ， \(u(x)\) 属于过 \(c\) 与 \(x\) 的直线，且满足关系 \(\Phi(x - c, u(x) - c) = \alpha\) 。滥用术语，我们说 \(u\) 是 \(\mathrm{L}\) 中的一个反演。
\end{enumerate}

\begin{enumerate}
\def\labelenumi{\alph{enumi})}
\item
  若 \(u\) 、 \(\nu\) 是两个具有同一极点 \(c\) 而幂分别为 \(\alpha\) 、 \(\beta\) 的反演，则 \(uv^{-1}\) 是这样一个位似在 L - \(\{c'\}\) 上的限制：该位似以 \(c\) 为中心、比为 \(\alpha\beta^{-1}\) （第 II 章，第 2 版，附录 II，习题 6）。
\item
  设 S 是一个含 c 的球面（习题 12）。证明 S - \{c\} 在以 c 为极点的反演下的像是一个垂直于连结 c 与 S 的中心的直线的超平面（滥用术语，称这个超平面为 S 在所述反演下的像）。逆命题成立。
\item
  设 S 是一个不含 c 的球面。证明：若 \(\varpi\) 是 c 关于 S 的幂（习题 12 b)），则 S 在以 c 为极点、以 \(\alpha\) 为幂的反演下的像，就是 S 在以 c 为中心、比为 \(\alpha/\varpi\) 的位似下的像。若 n = 2，且对每个 \(x \in S\) ，用 T（相应地，T′）表示 S（相应地， \(u(S)\) ）在点 x（相应地， \(u(x)\) ）处的切线，用 D 表示过 c、x 与 \(u(x)\) 的直线，证明有 \((\widehat{\mathrm{D},\mathrm{T}}) = (\widehat{\mathrm{T}^{\prime},\mathrm{D}})\) 。
\item
  设 \(S_{1}\) 、 \(S_{2}\) 是两个正交的球面（习题 12 d)）， \(S_{1}'\) 、 \(S_{2}'\) 是它们在一个以 c 为极点的反演下的像。若 c 既不属于 \(S_{1}\) 也不属于 \(S_{2}\) ，证明 \(S_{1}'\) 与 \(S_{2}'\) 是正交的球面。若 \(c \in S_{1}\) 且 \(c \notin S_{2}\) ，则 \(S_{1}'\) 是一个包含 \(S_{2}'\) 的中心的超平面。若 \(c \in S_{1} \cap S_{2}\) ，则 \(S_{1}'\) 与 \(S_{2}'\) 是互相垂直的超平面（§ 6，习题 22）。逆命题成立。
\item
  设 \(u\) 是以 \(c\) 为极点、以 \(\alpha = \rho^2 > 0\) 为幂的反演，C 是以 \(c\) 为中心、以 \(\rho\) 为半径的球面。若 \(x_1, x_2\) 是位于过 \(c\) 的直线上且异于 \(c\) 的两个不同点，则下列性质等价： \(\alpha\) \(x_1\) 与 \(x_2\) 在 \(u\) 下彼此互变； \(\beta\) \(x_1\) 与 \(x_2\) 关于 C 共轭（§ 6，习题 25）； \(\gamma\) 每个包含 \(x_1\) 与 \(x_2\) 的球面都与 C 正交。我们也说 \(u\) 是以 C 为球面的反演。
\end{enumerate}

¶ 14) 假设与记号同习题 12，在 L 中取定一个原点 0。设 E \(_{1}\) 是 L 与一个一维空间 Af \(_{1}\) 的向量空间直和；用 Q \(_{1}\) 表示 E \(_{1}\) 上的二次型，使得对 x ∈ L 与 η ∈ A，有

\[
\mathrm{Q} _ {1} (x + \eta f _ {1}) = \mathrm{Q} (x) + \eta^ {2},
\]

这个型是正的且非退化的；用 C 表示 E \(_{1}\) 中以 0 为中心、以 1 为半径的球面（关于 Q \(_{1}\) ）。

在欧几里得空间 E \(_{1}\) 中，设 s 是以 -f \(_{1}\) 为极点、以 2 为幂的反演（习题 13）；它在 L 上的限制 s \(_{0}\) 把 L 映为 C -\{-f \(_{1}\) \}；滥用术语，称 s \(_{0}\) （相应地， s \(_{0}^{-1}\) ）为以 -f \(_{1}\) 为视点、L 到 C 上（相应地，C 到 L 上）的球极投影。对 L 中每个以 c 为极点的反演 u， s \(_{0}\) us \(_{0}^{-1}\) 是 C 中集 \{s \(_{0}\) (c), -f \(_{1}\) \} 的补集的一个对合置换；令 u'(s \(_{0}\) (c)) = -f \(_{1}\) ， u'(-f \(_{1}\) ) = s \(_{0}\) (c)，把它开拓为 C 的一个对合置换 u′，并称 u′ 是 C 中的一个反演。类似地，对 L 中每个关于超平面的对称（变换）v， s \(_{0}\) vs \(_{0}^{-1}\) 是 C 中集 \{-f \(_{1}\) \} 的补集的一个对合置换；令 v'(-f \(_{1}\) ) = -f \(_{1}\) ，把它开拓为 C 的一个对合置换 v′，我们说 v′ 是 C 中的一个对称（变换）。由反演与对称（变换）生成的 C 的置换群称为 C 的共形群（滥用术语，或称 L 的共形群）。

\begin{enumerate}
\def\labelenumi{\alph{enumi})}
\item
  证明 C 的共形群由对称（变换） \(v'\) 与反演 \(u'\) 生成，它们对应于 L 中的反演 \(u\) ，其幂 \(>0\) 。（利用习题 13 a)，并注意在 L 中每个平移以及位似 \(x \to -x\) 都是关于超平面的对称（变换）之积）。
\item
  设 u 是 L 中幂 \textgreater{} 0 的一个反演；证明 C 中相应的反演 \(u'\) 是一个完全确定的变换 \(u_{1}'\) 在 C 上的限制，该变换或者是幂 \textgreater{} 0（在 \(E_{1}\) 中）且球面（习题 13 c)）与 C 正交的反演，或者是关于 \(E_{1}\) 中过 0 的一个超平面的对称（变换）（在 \(E_{1}\) 中考察与 u 有相同极点与相同幂的反演 \(u_{1}\) ）。对 C 中对应于 L 中关于超平面的对称（变换）v 的对称（变换） \(v'\) ，叙述相应的命题。
\item
  在向量空间 \(E_{2} = A \times E_{1}\) 中，考虑二次型 \(Q_{2}\) ，它使得对 \(\zeta \in A\) 、 \(y \in E_{1}\) 有
\end{enumerate}

\[
\mathrm{Q} _ {2} ((\zeta , y)) = \zeta^ {2} - \mathrm{Q} _ {1} (y),
\]

这个型非退化且符号差为 \((1, n + 1)\) 。把 \(E_{1}\) 与它在射影空间 \(\mathbf{P}(E_{2})\) 中的典范像等同（第 II 章，第 2 版，附录 III，第 4 小节）。证明（用 b) 的记号）： \(u'\) 也是某个射影线性映射 \(\overline{u}''\) 在 C 上的限制，后者是由一个变换 \(u''\) （正交群 \(\mathbf{O}(Q_{2})\) 的）过渡到商而得，这个变换是关于 \(E_{2}\) 中一个非迷向超平面的对称（变换）。对 \(\varphi'\) 叙述相应的命题。由此推出，L 的共形群同构于群 \(\mathbf{O}(Q_{2})\) 关于其中心的商群（利用命题 5 及 § 6 的习题 17 c)）。再由此推出，共形群的每个变换都是至多 \(n + 2\) 个 L 中的反演或对称（变换）的乘积（参见 § 6，习题 15 e)）。

\begin{enumerate}
\def\labelenumi{\alph{enumi})}
\setcounter{enumi}{3}
\tightlist
\item
  设 \(\Sigma\) 是以仿射空间 L 中的球面与超平面为元素的集。由 \(b\) ) 推出：存在一个从 \(\Sigma\) 到 \(\mathbf{P}(\mathrm{E}_2)\) 中一个集的补集上的双射，该集由 \(x \in \mathrm{E}_1\) 中满足 \(\mathrm{Q}_1(x) \leqslant 1\) 的那些元素组成，使得两个正交的球面对应于两个关于 C 共轭的点。
\end{enumerate}

\begin{enumerate}
\def\labelenumi{\arabic{enumi})}
\setcounter{enumi}{14}
\item
  把习题 12 至 14 的定义与结果推广到下述情形：A 是毕达哥拉斯域，且存在 \(\Phi\) 的标准正交基。
\item
  设 A 是一个交换域，V 是 A 上维数为奇数 \(2r + 1\) 的向量空间，F 是积空间 A × V， \(\Psi\) 是 F 上一个非退化的交错型；在维数为 \(2r + 1\) 的射影空间 P(F) 中，我们称集 \(C_{0}\) （即 F 的全迷向平面（关于 \(\Psi\) ）的典范像直线所成的集）为与 \(\Psi\) 相伴的（射影）线性丛。
\end{enumerate}

以下假定 A 是极大有序域；设 \(\Phi\) 是 V 上一个正的非退化对称型。设 D 是 F 中与 V（关于 \(\Psi\) ）正交的直线，它包含在 V 中；设 H 是 V 中与 D（关于 \(\Phi\) ）正交的超平面；在 F 中，H 是关于 \(\Psi\) 的非迷向子空间，因此它关于 \(\Psi\) 的正交补是一个与 H 互补且包含 D 的平面 P。在仿射空间 \(E = \{1\} \times V \subset F\) 中，我们也称由 F 的不包含于 V 的全迷向平面（关于 \(\Psi\) ）与 E 的交所成的集 C 为与 \(\Psi\) 相伴的（仿射）线性丛；直线 \(\Delta = P \cap E\) 称为线性丛 C 的轴（对于由度量型 \(\Phi\) 在 E 上定义的欧几里得空间结构而言）。

\begin{enumerate}
\def\labelenumi{\alph{enumi})}
\item
  证明：在 E 中，每个等于 \(\Delta\) 的一个方向向量的平移（第 II 章， \(2^{\text{e}}\) 版，附录 II，第 3 小节）都使 C 保持不变（参见 § 4，习题 6）。
\item
  设 \((e_i)_{0\leqslant i\leqslant 2r}\) 是 V 关于 \(\Phi\) 的一个标准正交基，使得 \(e_0\in \mathrm{D}\) ，并且有 \(\Psi (e_{2i - 1},e_{2i}) = \rho_{i} > 0\) （对 \(1\leqslant i\leqslant r,\Psi (e_j,e_k) = 0\) ，这是对不属于形状 \((e_{2i - 1},e_{2i})\) 或 \((e_{2i},e_{2i - 1})\) 的指标偶而言的）
\end{enumerate}

（§ 7，第 3 小节，命题 6）。在仿射空间 E 中取一个原点 \(a \in \Delta\) ，并令 \(\Psi(a, e_0) = \rho_0\) 。设 \(x\) 是属于由 \(e_{2i-1}\) 与 \(e_{2i}\) 所生成的平面的一个向量， \(y\) 是同一平面中满足 \(\Phi(x, y) = 0\) 、 \(\Phi(y, y) = 1\) 且 \(x \wedge y = \lambda e_{2i-1} \wedge e_{2i}\) （其中 \(\lambda > 0\) ）的向量。设 \(E_i\) 是 E 中由点 \(a\) 、 \(a + e_0\) 、 \(a + e_{2i-1}\) 、 \(a + e_{2i}\) 所张成的三维仿射线性簇， \(R_x\) 是 \(E_i\) 与（E 中的）由 C 的过点 \(a + x\) 的直线所生成的仿射超平面的交。证明：（关于 \(\Phi\) ）与平面 \(R_x\) （在 \(E_i\) 中）正交的直线的方向是直线 \(L_x\) （它属于平面 \(Ae_0 + Ay\) ），使得若 \(\theta = (\widehat{D, L_x})\) ，则有

\[
\operatorname{tg} \theta = \frac {\rho_ {i}}{\rho_ {0}} \sqrt {\Phi (x , x)}
\]

这里假定平面 \(Ae_{0} + Ay\) 的定向使 \(e_{0} \wedge y\) 为正。

¶ 17) a) 设 A 是一个交换域，J : \(\xi \to \overline{\xi}\) 是 A 的一个对合自同构，E 是 A 上的一个二维向量空间， \(\Phi\) 是 E 上一个非退化的埃尔米特半双线性型（非交错的）， \((e_1, e_2)\) 是 E 关于 \(\Phi\) 的一个正交基（ \(\S\) 6，第 1 小节，定理 1），使得

\[
\Phi (\xi_ {1} e _ {1} + \xi_ {2} e _ {2}, \eta_ {1} e _ {1} + \eta_ {2} e _ {2}) = \alpha \xi_ {1} \bar {\eta} _ {1} + \beta \xi_ {2} \bar {\eta} _ {2}
\]

（α 与 β 属于 J 的不变量所成的域 K）；令 γ = β/α。我们把点 \(\xi_{1}e_{1} + \xi_{2}e_{2} \in E\) 与环 B 的元素 \(\xi_{1} + \xi_{2}\rho\) 等同，B 由条件 \(\rho^{2} = -\gamma\) 与 \(\rho\xi = \bar{\xi}\rho\) （对 \(\xi \in A\) ）所定义（§ 3，习题 4 a)）。对每个 \(x = \xi_{1} + \xi_{2}\rho \in B\) ，令 \(\tilde{x} = \bar{\xi}_{1} - \xi_{2}\rho\) 及 N(x) = x \(\tilde{x}\) = \(\tilde{x}\) x，于是 \(x \to \tilde{x}\) 是代数 B（在 K 上）的一个对合反自同构，且有 \(\Phi(x, x) = \alpha N(x)\) 。证明：每个行列式等于其乘子的 Φ 的相似变换（也称为正相似变换）都可以以唯一的方式写成 \(x \to xy\) ，其中 y 是 E 的一个非迷向向量，且其乘子为 N(y)。

\begin{enumerate}
\def\labelenumi{\alph{enumi})}
\setcounter{enumi}{1}
\item
  先设 J 是恒等映射，于是 K = A。若 A 的特征数 ≠ 2，试用这个方法重新得出第 1 小节的结果。当 A 的特征数为 2 时，展开相应的理论（参见 § 4，习题 14；要区分两种情形，视 γ 在 A 中是否为平方元而定）。
\item
  设 \(J \neq 1\) ，于是 A 是 K 的一个可分二次扩张。这时 \(\Phi\) 必满足条件 (T)（§ 4，第 2 小节及习题 1）；若 \(\Phi\) 的指数为 0，则 B 是一个以 K 为中心的反射体（第 VIII 章，§ 11，习题 4）；若 \(\Phi\) 的指数为 1，则 B 同构于 \(\mathbf{M}_2(\mathrm{K})\) 。
\end{enumerate}

\(d)\) 设 \(\mathrm{J} \neq 1\) 且 A 的特征数 \(\neq 2\) ；这时有 \(\mathrm{A} = \mathrm{K}(\theta)\) ，其中 \(\theta^2 = -\delta \in \mathrm{K}\) 。若 S 是 \(\Phi\) 的正相似变换的群，H 是 E 中比为 \(\neq 0\) 且属于 K 的位似的群（一个同构于 \(\mathrm{K}^*\) 的群），证明群 S/H 同构于旋转群 \(\mathrm{O}^+(\mathrm{Q})\) ，其中 Q 是 K 上一个三维向量空间 F 上的非退化二次型，使得存在 F 的一个正交基 \((f_1, f_2, f_3)\) ，对此有

\[
\mathrm{Q} \left(\zeta_ {1} f _ {1} + \zeta_ {2} f _ {2} + \zeta_ {3} f _ {3}\right) = \gamma \zeta_ {1} ^ {2} + \delta \zeta_ {2} ^ {2} + \gamma \delta \zeta_ {3} ^ {2}
\]

（参见 § 9，习题 15）。

¶ 18) a) 设 A 是一个交换域， \(\xi \to \bar{\xi}\) 是 A 的一个对合自同构，E 是 A 上一个维数为偶数 \(2m\) 的向量空间， \(\Phi\) 是 E 上指数为 0 且满足条件 (T) 的非退化埃尔米特型， \(\Delta\) 是 \(\Phi\) 关于 E 的一个基的判别式。设 M(Φ) 是 \(\Phi\) 的相似变换的乘子所成的群（§ 4，习题 8）。证明 M(Φ) 是 A 中形如 \(\alpha\bar{\alpha} - (-1)^{m}\beta\bar{\beta}\Delta\) 的元素所成的乘法群的一个子群。（对 \(m\) 用归纳法论证，并利用习题 17 以及下面两个注记： \(1^{\circ}\) 若 \(u\) 是乘子为 \(\mu\) 的相似变换， \(x\) 是 E 的一个向量， \(y = u(x)\) 且 \(z = u(y)\) ，则存在一个酉变换 \(\nu\) 使 \(\nu(y) = y\) 且 \(\nu(z) = \mu x\) ； \(2^{\circ}\) 若 \(\alpha\) 、 \(\beta\) 、 \(\lambda\) 是 A 中三个 \(\neq 0\) 的元素，且存在 A 中的 \(a\) 、 \(b\) 、 \(c\) 、 \(d\) 满足条件 \(\lambda = a\bar{a} + \alpha c\bar{c}\) 与 \(\lambda = b\bar{b} + \beta dd\bar{d}\) ，则存在 A 中的 \(s\) 、 \(t\) 满足条件 \(\lambda = s\bar{s} - \alpha\beta t\bar{t}\) 。）

\begin{enumerate}
\def\labelenumi{\alph{enumi})}
\setcounter{enumi}{1}
\tightlist
\item
  设 K 是一个极大有序域， \(\mathrm{K}_{1} = \mathrm{K}((t_{1}))\) 是以 K 为系数、关于一个未定元 \(t_{1}\) 的形式幂级数域（第 IV 章，§ 5，第 7 小节）， \(\mathrm{A} = \mathrm{K}_{1}((t_{2}))\) 是关于第二个未定元 \(t_{2}\) 、以 \(K_{1}\) 为系数域的形式幂级数域。设 E 是 A 上的一个六维向量空间，Q 是 E 上一个非退化二次型，使得存在一个正交基 \((e_{i})\) ，对此有
\end{enumerate}

\[
\mathrm{Q} (\sum_ {i = 1} ^ {6} \xi_ {i} e _ {i}) = \xi_ {1} ^ {2} + \xi_ {2} ^ {2} + \xi_ {3} ^ {2} + \xi_ {4} ^ {2} + t _ {1} \xi_ {5} ^ {2} + t _ {2} \xi_ {6} ^ {2}.
\]

证明不存在 Q 的以 \(t_{1}t_{2}\) 为乘子的相似变换。

\backmatter
\chapter*{历史注记}
\phantomsection
\addcontentsline{toc}{chapter}{历史注记}

（注意——罗马数字指本注记之末所附的参考文献。）

以现代面貌出现的二次型理论，其历史很难追溯到十八 \(^{e}\) 世纪后半叶之前，而且，如我们将要看到的，它的发展首先是为了满足数论、分析与力学的需要。但是这一理论的基本概念实际上从“欧几里得”几何的开端起就已出现，并且构成其骨架。因此，若不至少概略地谈谈“初等几何”自古以来的发展，就无法追溯它的历史。当然，我们只能关注若干一般观念的演变，读者也不应指望在这里找到关于任何特定定理的历史的精确资料，对此我们只需引证专门的历史著作或教学著作（*）即已足够。不言而喻，当我们在下文谈到同一条定理在各种代数语言或几何语言中的各种可能解释时，绝不是说这些“翻译”在各个时代都像今天这样为人们所熟悉；恰恰相反，本注记的主要目的正是要表明，数学家们是如何极其渐进地意识到这些往往面貌迥异的问题之间的亲缘关系的；我们还要表明，他们在此过程中如何被引导着为古人留下的大量几何定理引入某种连贯性，并最终试图精确地划定“几何”一词应当理解成什么。

除了巴比伦人发现的求解二次方程的公式（(I)，pp. 183-189）之外，二次型理论主要概念的产生都必须在其几何伪装之下来考察。这些概念首先以距离的平方（在平面上或在三维空间中）的面貌出现，而相应的“正交性”概念则是借助于直角引入的，欧几里得把直角定义为平角的一半（《几何原本》第 I 卷，定义 10）；距离与直角这两个概念由勾股定理——欧几里得大厦的拱顶石（*）——联系在一起。角的概念在希腊数学中出现得似乎很早（希腊人很可能得之于巴比伦人，后者凭借长期的天文经验精于角的使用）。众所周知，在古典时期，只定义小于两直角的角（欧几里得的“定义”与他关于直线或平面给出的定义一样含糊而不可用）；定向的概念并未被提出，尽管欧几里得使用（既无公理又无定义）了一条直线把平面分成两个区域这一事实，并在必要处仔细地加以区分（**）。在这个阶段，平面旋转群的思想只是以很不完善的方式，通过（也是欧几里得不加解释地引入的）半直线非定向角的加法显现出来，而这种加法原则上也只在和至多等于两直角时才有定义（***）。至于三角学，它受到几何学家的轻视，而被留给测量员与天文学家；正是后者（阿里斯塔克斯、希帕恰斯，尤其是托勒密（V））建立了（平面的或球面的）直角三角形的边与角之间的基本关系，并编制了最早的数表（这些数表给出由角 \(\theta < \pi\) 在半径为 \(r\) 的圆上截得的弧的弦，换言之即数 \(2r \sin \frac{\theta}{2}\) ；更为便于使用的正弦的引入归功于中世纪的印度数学家）；在计算这些数表时，当时尚不知道的弧的加法公式被关于内接于圆的四边形的托勒密定理（可能上溯到希帕恰斯）的等价用法所代替（参见 Esp. vect. top.，第 V 章，§ 1，习题 5）。还应注意，欧几里得与海伦给出了与公式

\[
a ^ {2} = b ^ {2} + c ^ {2} - 2 b c \cos A
\]

等价的命题，它联系任意平面三角形的边与角；但由于缺乏直到十九 \(^{e}\) 世纪才出现的向量分析的思想，人们很难把这看作与度量型相伴的双线性型概念的首次露面。

位移（或运动，这两个概念的区分在古代——甚至在很久以后——并不清楚）已为欧几里得所知；但由于我们不知道的原因，他对使用它们似乎表现出明显的迟疑（例如在“三角形相等的情形”中，人们得到的印象是：他之所以使用位移的概念，只是因为未能表述出一条恰当的公理（(III bis)，t. I，pp. 225-227 及 249））；然而，正是借助于位移（绕轴旋转）的概念，他给出了旋转锥面与球面的定义（《几何原本》第 XI 卷，定义 14 与 18），阿基米德在定义旋转二次曲面时也是如此。但是，作用于整个空间的变换这一一般思想，在十八 \(^{e}\) 世纪末之前几乎不为数学思维所知（*）；

而在十七 \(^{e}\) 世纪之前，同样没有运动合成概念的踪迹，更不用说位移的合成了。这当然不意味着希腊人对图形的“正则性”与“对称性”不特别敏感——我们今天把这些性质与位移群的概念联系在一起；他们的正多边形理论、尤其是正多面体理论——他们全部数学中最出色的篇章之一——恰恰证明了相反的情形（*）。

最后，希腊数学在我们所关心的领域中的最后一项本质贡献是圆锥曲线的理论（至于二次曲面，希腊人只认识某些旋转二次曲面，而且除球面之外并未把研究推得很远）。这里值得注意的是：虽然希腊人从未有过解析几何基本原理的思想（主要是因为缺乏一种便于使用的代数），但他们在研究特殊“图形”时通常使用平面内相对于两条（甚至多于两条）轴的“纵标”（这些轴与图形密切相关，这正是他们的方法与费马和笛卡儿的方法根本不同之点之一，后者的轴是独立于所考虑的图形而固定的）。特别地，在与倍立方问题的联系中出现的最早一批（异于圆的）圆锥曲线，是由方程 \(y^{2} = ax\) 、 \(y = bx^{2}\) 、 \(xy = c\) 给出的曲线（门奈赫莫斯，欧多克索斯的学生， \(1v^{e}\) 世纪中叶）（**）；而在研究与这些曲线有关的问题时，最常用的是圆锥曲线的方程（通常相对于由一条直径与该直径和曲线的一个交点处的切线所成的两条斜轴）（而“焦点”性质只起很不显著的作用，与仅仅上溯到 \(x1x^{e}\) 世纪的学校传统可能使人相信的相反）。在这门庞大的理论中，我们在此尤须记住的是共轭直径的概念（阿基米德已经知道它），以及如今用作一点的极线之定义的那条性质，它由阿波罗尼奥斯（IV）在点位于圆锥曲线之外时给出（对他而言，极线就是从该点引出的各切线的切点所连成的直线）；从我们的观点看，这是关于异于度量型的二次型的“正交性”的两个例子，但这些概念与垂线的经典概念之间的联系在当时当然是不可能被想到的。

在笛卡儿与费马之前，几乎没有别的进展可报道；但从解析几何的一开始，二次型的代数理论就开始从它的几何外壳中显露出来：费马知道平面上一个二次方程表示一条圆锥曲线（(VII a)，pp. 100-102），并对二次曲面勾画了类似的想法（VII b）。随着十八 \(^{e}\) 世纪二维与三维解析几何的发展，（特别是联系圆锥曲线与二次曲面）出现了这一理论的两个中心问题：把二次型化归为平方和，以及寻求它关于度量型的“轴”。对圆锥曲线而言，这两个问题太初等，不会引起重要的代数进展；对任意多个变量的情形，第一个问题由拉格朗日在 1759 年联系多元函数的极大值问题解决（IX a）。但这个问题几乎立即被寻求轴的问题所掩盖，甚至在秩的不变性被表述之前就是如此（*）；至于惯性定律，它直到 1850 年前后才由雅可比（XIX b）与西尔维斯特（XX）发现，前者用与今天所用相同的推理证明了它，后者则只是把它当作几乎显然的事实加以陈述（**）。

把二次曲面化归到其轴的问题已呈现出比圆锥曲线的相应问题大得多的代数困难；最早着手研究它的欧拉不能证明特征值的实性，他在一段缺乏证明价值的概述之后承认了它（(VIII a)，pp. 379-392）（**）。即使这一点在 1800 年前后得到正确确立，仍须等到柯西才对任意变量个数 \(n\) 的型证明了相应的定理（XIV b）。也是柯西，在大约同一时间，证明了给出特征值的特征方程在直角轴的任何改变下不变（(XIV a)，p. 252）（****）；但对 \(n = 2\) 或 \(n = 3\) ，由于借助相应的圆锥曲线或二次曲面的轴对特征值所作的几何解释，这种不变性在直观上是“显然”的。此外，在关于这个课题的研究过程中，特征值的初等对称函数也自然而然地出现（带有种种几何解释，特别涉及阿波罗尼奥斯关于共轭直径的定理），特别是判别式，它（就 n = 2 而言，早在与二次方程理论的联系中便已为人所知）在 n = 3 时由欧拉首次给出（(VIII a) p. 382）；欧拉是在对二次曲面分类时遇到它的（在表达一个二次曲面没有任何无穷远点的条件时），而未提到它在直角轴改变下的不变性。但稍后，随着整系数二次型算术理论的开端，拉格朗日（就 n = 2）注意到判别式在线性但非正交的变量替换下的不变性的一个特例（(IX b)，p. 699），而高斯对 n = 3 建立了判别式在任何线性变换下的“协变性”（(XI a)，pp. 301-302）（*）。在柯西与比内证明了行列式乘法的一般公式之后，把高斯的公式推广到任意多个变量的情形便是直接的了；正是它在大约 1845 年给出了不变量一般理论的最初推动。

在希腊人那里取代位移理论的两个概念——即运动的概念与图形的“对称”的概念——之外，十七 \(^{e}\) 与十八 \(^{e}\) 世纪又加上了第三个，即直角轴变换的问题，它与这个理论实质上等价。欧拉有好几部著作致力于这个问题，他尤其力求为轴变换公式得到便于使用的参数表示。我们知道力学后来对他为此目的在 \(n = 3\) 时引入的三个角作了怎样的使用（(VIII a)，pp. 371-378）。但他并未止步于此；他在 1770 年考虑了对任意 \(n\) 的正交变换的一般问题，注意到通过引入 \(n(n - 1)/2\) 个角作为参数便可达到目的，最后，对 \(n = 3\) 与 \(n = 4\) 给出了旋转的有理表示（分别用 4 个齐次参数与 8 个由一个关系式联系起来的齐次参数表示），它们正是后来借助四元数理论得到的那些表示（参见 § 9，习题 15 与 16），而他并未指出其来源（VIII c）（**）。

另一方面，欧拉还指出了如何解析地寻求平面图形的“对称”，正是在这一场合，他实质上证明了平面位移是一个旋转，或一个平移，或一个平移加一个对称（(VIII a)，pp. 197-199）。当时力学的兴起还导致对位移的一般研究；但起初只有与连续运动相切的“无穷小”位移才是讨论的对象：在托里拆利、罗贝瓦尔与笛卡儿关于平面运动的合成与转动瞬时中心的研究中，出现的显然只有这种位移（参见第 IV 卷的历史注记，第 I-II-III 章）。转动瞬时中心由约翰·伯努利以一般的方式定义；达朗贝尔在 1749 年、欧拉在次年推广了这一概念，证明了使一点保持不动的运动存在转动的瞬时轴。有限位移的类似定理直到 1775 年才由欧拉陈述（VIII d），在这篇论文中他还同时发现旋转的行列式等于 1；次年，他证明了平面相似变换存在不动点（VIII e）。但必须等到夏尔自 1830 年起的工作（XV a），才最终得到关于有限位移与无穷小位移的一个连贯的理论。
